%=====================================================================
\section{Numerische Details und Robustheit}\label{sec:numerics}
%=====================================================================
\begin{itemize}
  \item Durchgehend \verb|std::complex<double>| bzw.\ \verb|cuDoubleComplex|; keine einfache Genauigkeit
        im Kern. (Gemischte Genauigkeit nur optional für LU mit iterativer Verfeinerung, M10.)
  \item Matrizen nie explizit invertieren, wenn ein Gleichungssystem genügt (LU + Lösen).
        Ausnahme: $\Li$, $\tp{C_\varepsilon}^{-1}$, $Q_{\varepsilon,yy}^{-1}$, die mehrfach gebraucht werden
        (dann LU speichern und Produkte per Lösen bilden).
  \item Eigenwertklassifikation mit Toleranzen relativ zu $\max|\lambda|$; entartete Eigenwerte
        (Symmetrie, Normaleinfall) sind erlaubt -- die Eigenvektoren spannen den Raum auf, die
        S-Matrix ist davon unabhängig.
  \item Rayleigh-Anomalien: \eqref{eq:beta_branch}; zusätzlich Option, die Wellenlänge um
        $10^{-12}$ relativ zu verschieben.
  \item Diagnosen je Rechnung: Energiebilanz $\sum e^{(r)}+\sum e^{(t)}$ (verlustfrei $=1$), Konditionszahlen
        ($T_{11}$, $G$, $Q_{\varepsilon,yy}$), Anzahl RK-Schritte, Klassifikationsreserve.
  \item Reproduzierbarkeit: deterministische Reihenfolgen, keine Zufallszahlen im Kern.
  \item \textbf{Eigen-Falle:} \verb|dot()| konjugiert das erste Argument, \verb|cross()| liefert für komplexe
        Vektoren das konjugierte Kreuzprodukt. Bilineare Formen ($\xi=\mathbf N^T\overline{\overline\varepsilon}\mathbf N$,
        $\mathbf v\times\mathbf T$, $\mathbf E\cdot\hat{\mathbf s}$) explizit ausschreiben (Befund M5).
  \item Symmetrien (optional, später): spiegelsymmetrische Strukturen bei Normaleinfall erlauben
        Halbierung der Ordnungen -- \emph{nicht} in der ersten Version.
\end{itemize}

%=====================================================================
\section{CUDA-Backend (RTX 6000 Ada)}\label{sec:cuda}
%=====================================================================

\subsection{Hardwareeigenschaften und Konsequenzen}
RTX 6000 Ada: Compute Capability 8.9 (\verb|sm_89|), 48\,GB GDDR6. Die FP64-Rate beträgt bei dieser
Architektur nur einen kleinen Bruchteil der FP32-Rate (Größenordnung 1/64; vor M10 mit einem
ZGEMM-Benchmark messen). Konsequenzen:
\begin{itemize}[nosep]
  \item Die GPU lohnt sich vor allem bei großen dichten Matrizen (2D, $n\gtrsim1000$) und bei vielen
        unabhängigen Rechnungen (Wellenlängen-/Winkel-Scans, gebatcht).
  \item Für 1D-Probleme ($n\le$ einige Hundert) ist der CPU-Pfad meist gleich schnell oder schneller;
        die Backend-Auswahl erfolgt automatisch nach Problemgröße (Schwelle konfigurierbar).
  \item Der CPU-Pfad bleibt Referenz für alle Tests (GPU-Ergebnisse müssen auf $10^{-10}$ relativ übereinstimmen).
\end{itemize}

\subsection{Bibliotheken}
\begin{center}
\small
\begin{tabular}{p{2.6cm}p{6.3cm}p{5.3cm}}
\toprule
Operation & GPU & CPU \\
\midrule
GEMM & cuBLAS \verb|cublasZgemm| (+ \verb|StridedBatched|) & Eigen / BLAS (OpenBLAS, FlexiBLAS) \\
LU / Lösen & cuSOLVER \verb|cusolverDnXgetrf|, \verb|Xgetrs| (64-Bit-API) & Eigen \verb|PartialPivLU| / LAPACK \verb|zgetrf| \\
Eigenwerte (allgemein) & cuSOLVER \verb|cusolverDnXgeev| (ab CUDA 12.6 Update 2; hybrid CPU/GPU, gepinnter Host-Speicher) & Eigen \verb|ComplexEigenSolver| / LAPACK \verb|zgeev| \\
FFT & cuFFT (Z2Z) & FFTW3 oder PocketFFT (header-only) \\
\bottomrule
\end{tabular}
\end{center}
Eigene Kernel: Aufbau (block-)Toeplitz-Matrizen aus Koeffizientengittern, Diagonalskalierungen
($\alpha X$, $X\gamma$), punktweise Produkte $C_\varepsilon(x)\mathbf F(x)$, Feldsynthese-Nachbearbeitung.
Zielarchitektur: \verb|CMAKE_CUDA_ARCHITECTURES=89|.
Speicherverwaltung über RAII-Klassen (\verb|DeviceBuffer<T>|), Streams je Wellenlänge,
Fehlerprüfung per Makro nach jedem Aufruf.

\subsection{Schnittstelle}
Der Physikcode ruft nur die abstrakte Schnittstelle \verb|fg::linalg::Backend| (GEMM, LU, Solve, Eig,
FFT, Speichertransfers) auf. Implementierungen: \verb|CpuBackend| (Pflicht, zuerst), \verb|CudaBackend|
(M10). Kein CUDA-Code außerhalb von \verb|src/cuda/| und \verb|src/linalg/cuda_backend.*|.

\subsection{Umsetzung (M10, v1.8)}
\begin{itemize}[nosep]
  \item \verb|CudaBackend| (\verb|src/cuda/cuda_backend.cpp|) nutzt nur die Host-API: \verb|cublasZgemm|,
        \verb|cusolverDnXgetrf|/\verb|Xgetrs| (64 Bit), \verb|cusolverDnXgeev| (rechte Eigenvektoren, auf
        Einheitsnorm normiert wie Eigen; Host-Workspace gepinnt). RAII für Device-/Pinned-Speicher, je Backend
        ein Stream, Fehlerprüfung per Makro. Ohne eigene Kernel ist kein \verb|nvcc| nötig
        (\verb|find_package(CUDAToolkit 12.6)|, \verb|CUDA::cudart|, \verb|cublas|, \verb|cusolver|); eigene
        Kernel (Toeplitz-Aufbau, Synthese) und cuFFT folgen bei Bedarf (M11).
  \item Dispatch: Die $O(n^3)$-Stellen (RCWA-Eigenprobleme, stabile Schicht-S-Matrix, Sternprodukt, $T\to S$,
        RK-Schrittpropagator) fragen \verb|linalg::offload(n)|; ist es falsch, läuft der unveränderte
        Eigen-Code (CPU-Referenz bitgleich), sonst \verb|linalg::device()|. Einstellung thread-lokal über
        \verb|SolverOptions::backend| (\verb|auto|, \verb|cpu|, \verb|cuda|) und \verb|backendThreshold| (1000).
  \item Batching über Wellenlängen: \verb|solveScan| verteilt die Scanpunkte auf OpenMP-Threads, jeder mit
        eigenem Löser und (GPU) eigenem Backend/Stream; Ergebnisse identisch mit sequentiellen Rechnungen.
  \item Ohne GPU (Sandbox) werden die GPU-Tests übersprungen; Bau und Linken gegen CUDA 12.8 sind geprüft,
        die Abnahme GPU $=$ CPU ($10^{-10}$) und der Benchmark-Bericht (\verb|docs/benchmarks.md|) erfolgen auf
        dem Zielsystem.
\end{itemize}

%=====================================================================
\section{Software-Architektur}\label{sec:arch}
%=====================================================================

\subsection{Verzeichnisse}
{\footnotesize\begin{verbatim}
include/fg/core      Typen, Konstanten, Fehler, IndexMap
include/fg/fourier   Koeffizienten (Intervall/Rechteck/Polygon/FFT), Toeplitz
include/fg/geometry  Material, Schichttypen, Stapel, Normalenfelder, Scheibenzerlegung
include/fg/physics   Einfall, Q-Matrizen (FFF), Propagationsmatrix M, Modenbasen, Effizienzen
include/fg/smatrix   S-Matrix, Sternprodukt, T->S, Gesamtlösung, interne Amplituden
include/fg/solvers   Butcher-Tableaus, RK-Löser, RCWA-Löser, gemeinsame Schnittstelle
include/fg/field     Nahfeld: Schnitte, Rekonstruktion, Synthese
include/fg/linalg    Backend-Schnittstelle, CPU- und CUDA-Backend
include/fg/io        JSON-Konfiguration, Ergebnis- und Feldausgabe
src/...              Implementierungen (gleiche Gliederung), src/cuda/ Kernel
apps/fgsolve         Kommandozeilenprogramm
tests/unit           GoogleTest, je Modul
tests/validation     Referenzfälle (Abschnitt Validierung), CTest-Label "validation"
examples/            JSON-Konfigurationen der Referenzfälle
scripts/             Setup AlmaLinux 10, Plot-Skripte, Doku-Build
docs/                dieser Plan, Architektur, Roadmap, Validierung, Konventionen
\end{verbatim}}

\subsection{Zentrale Typen (Skizze)}
{\footnotesize\begin{verbatim}
using cplx = std::complex<double>;  using MatC = Eigen::MatrixXcd;  using VecC = Eigen::VectorXcd;
struct Material { bool isotropic; cplx eps; Eigen::Matrix3cd epsTensor; };   // mu = 1
struct Incidence { double lambda_um, phi_deg, psi_deg, delta_deg; };
enum class Dim { D1, D2 };   enum class Factorization { Laurent, LiLamellar, FFF_iso, FFF_A, FFF_B };
struct Truncation { int Nx, Nz; };                      // Nz = 0 im 1D-Fall
class IndexMap;                                         // (n,m) <-> r
struct LayerSpec { /* variant: Homogeneous, Lamellar1D, Profile1D, Polygon2D, Profile2D */ };
struct Stack { Material substrate, superstrate; std::vector<LayerSpec> layers; double dx, dz; };
struct QBlocks { MatC b[3][3]; };                       // Q_eps,ij
MatC buildPropagationMatrix(const QBlocks&, const VecC& alpha, const VecC& gamma); // eq:Mgeneral
struct ModeBasis { MatC Psi, PsiInv; VecC lambda; };    // homogene Medien
struct SMatrix { MatC S11, S12, S21, S22; };            // Konvention eq:Sdef
SMatrix star(const SMatrix& lower, const SMatrix& upper);
SMatrix fromTransfer(const MatC& T);
class Solver { virtual SMatrix sliceSMatrix(...) = 0; };  // RKSolver, RCWASolver
struct Result { std::vector<OrderEfficiency> R, T; double absorption, energyBalance; Diagnostics diag; };
struct NearFieldRequest { double y_horizontal; double z_vertical; GridSpec grid; Quantities q; };
\end{verbatim}}

\subsection{Konfiguration (JSON, Kurzfassung; Details in \texttt{docs/config\_format.md})}
{\footnotesize\begin{verbatim}
{ "dimension": "1D", "solver": "rk" | "rcwa",
  "truncation": {"Nx": 15, "Nz": 0}, "factorization": "auto",
  "period_um": {"x": 0.5, "z": 0.0},
  "incidence": {"lambda_um": 0.6328, "phi_deg": 30, "psi_deg": 0, "delta_deg": 90},
  "superstrate": {"n": 1.0}, "substrate": {"n": [1.38, 7.62]},
  "layers": [{"type": "profile1d", "shape": "sinus", "depth_um": 0.2,
              "below": "substrate", "above": "superstrate", "slices": 40}],
  "rk": {"tableau": "gauss2", "tol": 1e-8}, "rcwa": {"slices": 60},
  "nearfield": {"horizontal_y_um": 0.1, "vertical_z_um": 0.0, "nx": 400, "ny": 300,
                "quantities": ["E","H","absE2"], "reconstruction": "fff"},
  "output": {"dir": "out/al_sinus"} }
\end{verbatim}}

\subsection{Code-Regeln}
C++20, \verb|clang-format| (Datei im Repo), keine rohen \verb|new|/\verb|delete|, Fehler über
\verb|fg::Error| (Ausnahmen) mit Kontext, keine globalen Zustände, Einheiten nur an der IO-Grenze
umrechnen, jede öffentliche Funktion mit Kommentar inkl.\ Formelreferenz auf diesen Plan
(z.\,B.\ \verb|// Plan (7.3)|).

%=====================================================================
\section{Validierung und Tests}\label{sec:validation}
%=====================================================================

\subsection{Unit-Tests (Pflicht, je Modul)}
\begin{enumerate}[nosep]
  \item IndexMap: Bijektion $(n,m)\leftrightarrow r$ für alle Paare; 1D-Spezialfall.
  \item Fourierkoeffizienten: Intervall vs.\ numerische Quadratur; Polygon \eqref{eq:polygon} vs.\ Rechteck;
        Summe der Koeffizienten zweier komplementärer Gebiete $=\delta_{p0}$.
  \item Toeplitz: Laurent-Regel für glatte Funktionen (Inneres exakt); Beispiel Buch IV.1
        (inverse Regel bei $N=0$ exakt $=1$, Laurent $1{,}5625$).
  \item Butcher-Tableaus: Ordnungsbedingungen; $P_h$ vs.\ \verb|expm| (Padé-Ordnung).
  \item Propagationsmatrix: homogenes isotropes Medium $\Rightarrow$ $M\Psi=\Psi\Lambda$ mit \eqref{eq:Psi_iso};
        Vergleich der 16 Blöcke mit direkter Elimination aus \eqref{eq:mx1}--\eqref{eq:mx3} (Zufalls-$Q$).
  \item S-Matrix: Sternprodukt vs.\ T-Produkt; neutrales Element; Assoziativität.
  \item $Q_\varepsilon$: Grenzfälle aus Abschnitt~\ref{sec:fff}; $A_\varepsilon C_\varepsilon=\mathbb 1$ für Zufallstensoren;
        Variante B isotrop $\approx$ \eqref{eq:Qiso}.
  \item Modenbasis anisotrop: isotroper Tensor reproduziert \eqref{eq:Psi_iso}; uniaxial mit optischer Achse
        $\parallel y$: $\lambda=\pm i\sqrt{\varepsilon_o-\kappa^2}$ (TE-artig) und $\pm i\sqrt{\varepsilon_o(1-\kappa^2/\varepsilon_e)}$.
\end{enumerate}

\subsection{Referenzfälle (Validierung)}
\begin{center}
\small
\begin{tabular}{p{0.6cm}p{6.6cm}p{6.6cm}}
\toprule
ID & Fall & Referenz / Kriterium \\
\midrule
V1 & ebene Grenzfläche, alle Winkel, TE/TM, beliebiger Azimut & Fresnel analytisch, $10^{-12}$ \\
V2 & homogene Schicht auf Substrat (Airy), auch absorbierend & analytisch, $10^{-12}$ \\
V3 & Lamellengitter dielektrisch ($\nu=1{,}5$ bzw.\ $2{,}5$), TE/TM/konisch & Energie $=1\pm10^{-10}$;
     RK = RCWA auf RK-Toleranz; Reziprozität nur bis auf Abschneidung (die Ordnungsmengen beider
     Einfallsrichtungen sind um eine Ordnung verschoben): Abweichung muss mit $N$ fallen \\
V4 & Al-Sinusgitter: $d=0{,}5\,\mu$m, $g=\tfrac a2(1+\cos Kx)$, $a=0{,}2\,\mu$m, $\lambda=0{,}6328\,\mu$m,
     $\Phi=30^\circ$, $\nu=1{,}38+7{,}62i$, Ordnung $-1$ & TE $0{,}3235$, TM $0{,}8696$ (intern, C-Methode;
     FFF-Treppe 300 Scheiben $0{,}8687$); Toleranz $2\cdot10^{-3}$ \\
V5 & tiefes Glas-Sinusgitter, $d=0{,}5$, $a=1{,}0\,\mu$m, $\nu=1{,}5$, $30^\circ$, Transmission $-1$ & TE $0{,}820$
     (intern, 120 Scheiben); Buch Abb.\ III.4 $\approx0{,}82$ \\
V6 & Dielektrisches Sinusgitter $\nu=2{,}5$, $a/d=1$, $d=1\,\mu$m, $30^\circ$ (Buch Abb.\ IV.3), TM &
     FFF-Fehler $10^{-4}$ bei $N\approx8$ (qualitativ: FFF $\gg$ klassisch) \\
V7 & Kobalt-Sinusgitter anisotrop (Buch Kap.\ VIII): $d=0{,}6$, $a=0{,}5\,\mu$m, $30^\circ$, TM-Einfall,
     $\varepsilon_{ii}=-8{,}19+16{,}38i$, $\varepsilon_{xz}=-\varepsilon_{zx}=-0{,}495-0{,}106i$ &
     intern ($N=20$, 100 Scheiben): TM$_0$ $0{,}0567$, TM$_{-1}$ $0{,}1435$, TE$_0$ $1{,}3\cdot10^{-7}$,
     TE$_{-1}$ $2{,}37\cdot10^{-6}$; Buch $\approx0{,}055/0{,}148/\approx10^{-7}/2{,}6\cdot10^{-6}$; Toleranz 3\,\%.
     M5: $0{,}05673/0{,}14355/1{,}275\cdot10^{-7}/2{,}372\cdot10^{-6}$; RK ($N=8$) = RCWA (400 Scheiben) auf $10^{-4}$ \\
V7b & einachsiges Substrat, Achse $\parallel y$, $\varepsilon_o=2{,}4$, $\varepsilon_e=3{,}1$ & analytisch ($r_s$, $r_p$ s.\ Text), $10^{-12}$ \\
V7c & verlustfreie biaxiale gedrehte Schicht, konisch & Energie $10^{-12}$; homogener Pfad = Lamellenpfad (RCWA $10^{-10}$, RK $10^{-8}$) \\
V8 & Anisotroper Code mit isotropem Tensor & identisch mit isotropem Code ($10^{-10}$ bei gleicher Faktorisierung).
     2D (M8): Laurent mit Tensor = isotrop (Polygon, Eierkarton); \verb|FFF_B| auf $z$-invariantem Streifen = Li;
     anisotrope $z$-invariante Strukturen = 1D-anisotrop ($10^{-10}$/$10^{-12}$); Spiegeltest $x\leftrightarrow z$
     mit gespiegeltem Tensor ($10^{-9}$). Konvergenz: Co-Säulen \verb|FFF_B| $T=0{,}548\ldots0{,}551$ ($N=2\ldots7$),
     Laurent $0{,}42\ldots0{,}49$; anisotroper Eierkarton: A, B, Laurent konsistent \\
V9 & 2D-Code mit $z$-invarianter Struktur ($g(x,z)=g(x)$, $N_z=0$) & identisch mit 1D-konisch.
     M7: Lamellen, \texttt{Polygon2D}-Streifen (NF-Tensor, \verb|FFF_sym|/\verb|FFF_iso| = Li),
     \verb|sinus_x| und 1D-Profil im 2D-Stapel, RK und RCWA: $10^{-12}$ ($N_z=0$); RCWA mit $N_z>0$: Ordnungen
     $(n,0)$ $10^{-10}$, übrige Ordnungen leistungsfrei; RK mit $N_z>0$ nur auf RK-Toleranz (Schrittweitensteuerung
     sieht die entkoppelten Blöcke) \\
V10 & Lochgitter (Buch IX.1.7): Ag-Film $0{,}2\,\mu$m auf Glas, quadratische Löcher $0{,}25\,\mu$m,
      Periode $0{,}9\,\mu$m, $\nu_{\text{Ag}}=0{,}1+8{,}94i$, Normaleinfall & $T(\lambda=1{,}388\,\mu$m$)\approx7{,}5\,\%$
      (Buch); Plasmonpol der ebenen Grenzfläche bei $1{,}369406\,\mu$m. M7 (RCWA, \verb|FFF_sym|, NF-Tensor):
      Peak stets oberhalb des Pols; $\lambda_{\text{Peak}}=1{,}383/1{,}386/1{,}391/1{,}404/1{,}3876/1{,}3889\,\mu$m und
      $T_{\max}=4{,}1/3{,}4/10{,}3/5{,}3/4{,}2/6{,}2\,\%$ für $N=5\ldots10$: ab $N=9$ Peak bei der Buch-Wellenlänge,
      Höhe nähert sich 7,5\,\%; Tensor- und Vektorvariante von Alg.~NF gleichwertig. Offen: $N=11\ldots15$
      (Zielsystem). Test: qualitativ, Polarisationsunabhängigkeit $10^{-9}$ \\
V11 & Fig.\ VI.10 (Al, $40^\circ$, $\nu=1{,}3+7{,}6i$): Li-Lamellentreppe vs.\ FFF & qualitativ: Treppe mit
      Li-Formel konvergiert langsam, FFF schnell \\
V13 & Gold-Sinus TM: $d=1$, $a=0{,}25\,\mu$m, $\nu=0{,}2+3{,}4i$, $20^\circ$ &
      C-Methode: $e_{-1}=0{,}26614$, $e_0=0{,}02362$; \verb|FFF_rig| $10^{-3}$, \verb|FFF_sym| $3\cdot10^{-3}$;
      Buchvariante versagt (dokumentiert) \\
V14 & Gold-Dreieck TM (Spitze bei $0{,}7\,\mu$m, sonst wie V13) & offen: nur physikalische Grenzen
      ($\sum e<1$); Laurent bei $N\le64$ noch nicht konvergiert \\
V12 & Nahfeld: Airy-Schicht analytisch; Stetigkeit an allen Element-/Scheibengrenzen (RK, RCWA, TE/TM);
      Flusskonstanz (verlustfreie Lamellen); RK-Feld = RCWA-Feld; $D_x$-Stetigkeit direkt vs.\ FFF &
      $10^{-12}$; $10^{-6}$; $10^{-12}$/$10^{-8}$; $10^{-7}$; FFF $<0{,}1$, direkt $>1$.
      2D (M9): $z$-invariante Strukturen = 1D-Nahfeld ($10^{-9}$, Karten = Linie $\cdot e^{i\gamma_0(z-z_0)}$),
      Stetigkeit (RK, RCWA) $10^{-6}$, Flusskonstanz $10^{-12}$, RK = RCWA $10^{-7}$, $D_x$-Sprung an der
      Säulenwand direkt 293\,\%, FFF 5\,\% / 3\,\% (anisotrop) \\
\bottomrule
\end{tabular}
\end{center}
\enquote{Intern} = in der Prüfung der Buchkapitel mit unabhängigen Python-Implementierungen berechnet
(nicht aus der Literatur); bei Abweichungen zuerst die eigene Implementierung, dann diese Werte prüfen.
Zusätzlich für jeden Referenzfall eine Konvergenzstudie in $N$ und (RCWA) in der Scheibenzahl als
Testausgabe (CSV) ablegen.

%=====================================================================
\section{Meilensteine}\label{sec:milestones}
%=====================================================================
\begin{longtable}{p{1.0cm}p{6.6cm}p{6.6cm}}
\toprule
M & Inhalt & Abnahme \\
\midrule
\endhead
M0 & Build-System, Abhängigkeiten, CI (CPU), Grundgerüst, Doku & \verb|cmake --build| + \verb|ctest| grün auf AlmaLinux 10 und in CI \\
M1 & IndexMap, Fourierkoeffizienten, Toeplitz, Butcher, S-Matrix-Kern, isotrope Modenbasis, $M$-Aufbau & Unit-Tests 1--6, 8 (isotrop) \\
M2 & 1D in-plane isotrop: homogene Schichten, Lamellen (RCWA, Li und FFF\_iso), Gesamtlösung, Effizienzen & V1, V2, V3 (TE/TM) \\
M3 & 1D Profile: Nullstellen-Schnitte, Normalenfeld, FFF\_iso; RK-Löser (Gauß 2/3, Schrittsteuerung); RCWA-Treppe mit FFF & V4, V5, V6; RK = RCWA auf Lamellen \\
M4 & 1D konisch ($4n_f$): volle $M$, Polarisation, s/p-Ausgabe & V1/V2 konisch, V3 konisch, Reziprozität \\
M5 & 1D anisotrop: Tensoren, FFF\_A, FFF\_B, anisotrope Substrate (4×4-Moden, Klassifikation) & V7, V8, Unit-Test 7 \\
M6 & Nahfeld 1D: interne Amplituden, RCWA/RK-Rekonstruktion, direkt/FFF, Schnitte, Plot-Skript & V12 (1D) \\
M7 & 2D isotrop: IndexMap 2D, Polygon-Koeffizienten, Normalenfeld (Alg.\ NF), RCWA; RK für Profile2D & V9, V10
  (erledigt v0.6; V10 nur qualitativ, quantitativ offen) \\
M8 & 2D anisotrop & V8 (2D), Konvergenzstudie (erledigt v0.7) \\
M9 & Nahfeld 2D (horizontal $x$--$z$, vertikal $x$--$y$) & V12 (2D) (erledigt v0.8) \\
M10 & CUDA-Backend (cuBLAS, cuSOLVER getrf/geev, cuFFT), Batching über Wellenlängen, Benchmarks & GPU = CPU auf $10^{-10}$; Benchmark-Bericht
  (v0.9: implementiert, gebaut und gelinkt; GPU-Abnahme auf dem Zielsystem offen) \\
M11 & Parameter-Scans, Performance, Doku für Anwender & Beispielskripte, Laufzeittabelle
  (erledigt v0.10: Scans $\lambda\times\varphi\times\psi\times\delta$, Sweeps per JSON-Pointer, paralleler Scan,
  \verb|plot_efficiencies.py|, \verb|docs/user_guide.md|) \\
\bottomrule
\end{longtable}

%=====================================================================
\section{Buchfehler mit Relevanz für die Implementierung}\label{sec:errata}
%=====================================================================
\begin{center}
\small
\begin{tabular}{p{2.3cm}p{5.4cm}p{6.6cm}}
\toprule
Stelle & Buch & korrekt (in diesem Plan verwendet) \\
\midrule
(II.19)/(II.20) & $\beta_0^{(1)}$ im einfallenden Term & $\beta_0^{(2)}$ \\
TM-Matrix S.\,46 & Indizes $(m,n)$ vertauscht & $M_{12,nm}=(k^2)_{n-m}$ usw. \\
(III.24)/(III.26) & $C$ und $C'$ vertauscht & Grenzflächen über $\Psi_{\text{ref}}^{-1}\Psi_X$ \\
S.\,56 & $A'^{(1)}=A^{(1)}$ & nur für $\nu_1=\nu_M$; Grenzfläche Substrat$\to$ref nötig \\
(III.29) & $T_{12}$, $T_{21}$ vertauscht & -- (Sternprodukt verwenden) \\
(III.45) & $S_{21}$ statt $S_{22}$ & \eqref{eq:star} \\
(IV.39)/(IV.40) & $\delta_{n,p}$ fehlt & Startwerte als Einheitsmatrix (Abschn.\ RK) \\
konisch & $\beta$ ohne $\gamma_0$ & $\beta=\sqrt{\varepsilon-\alpha^2-\gamma^2}$ \\
S.\,102 (Text) & Spaltenreihenfolge von $C_\varepsilon$ & \eqref{eq:Ceps} \\
S.\,108 & $C_\varepsilon$ $y$-unabhängig & hängt über $\varepsilon(y)$ von $y$ ab \\
(VI.13) & $L$ in allen Blöcken & nur Diagonalblöcke (schräge Lamellen) \\
(VII.72) & Vorfaktor $\sqrt{\mu_0/\varepsilon_0}/(ik_j)$ & $1/(i\omega\varepsilon_j)$ (nicht benötigt) \\
(VIII.2) & Vorzeichen $C_{23}$ & \eqref{eq:C1d} \\
(VIII.3) & $M_{32}$ ohne $-i\alpha^2/(\omega\mu_0)$ & \eqref{eq:Mmu1b} \\
(IX.1) & $\beta_0=k_0\cos\Psi$ & $\cos\Phi$, \eqref{eq:incidence} \\
(IX.14) & Division durch $N+1$ & \eqref{eq:indexmap} \\
(IX.16)--(IX.19) & Indizes $m(s)$, $\alpha_{m(r)}$ & alle zur Komponente $r$, $\alpha_{n(r)}$ \\
(IX.20) & Transmission ohne $k_2\to k_1$ & \eqref{eq:eff_closed} \\
Anh.\ IV & 2D-Blockgröße $(4N+2)^2$ & $2(2N+1)^2$ \\
(IV.32) & Vertauschung $\tp{N_i}\leftrightarrow\Li$ (Näherung) & versagt bei schwach absorbierenden Metallen;
  \eqref{eq:Qrig}/\eqref{eq:Qsym} \\
\bottomrule
\end{tabular}
\end{center}

%=====================================================================
\section*{Literatur}
\addcontentsline{toc}{section}{Literatur}
\begin{enumerate}[label={[\arabic*]},nosep,leftmargin=2em]
  \item M.~Nevière, E.~Popov, \emph{Light Propagation in Periodic Media: Differential Theory and Design}, Marcel Dekker, 2003.
  \item L.~Li, \enquote{Use of Fourier series in the analysis of discontinuous periodic structures}, JOSA A 13, 1870 (1996).
  \item L.~Li, \enquote{Formulation and comparison of two recursive matrix algorithms for modeling layered diffraction gratings}, JOSA A 13, 1024 (1996).
  \item L.~Li, \enquote{New formulation of the Fourier modal method for crossed surface-relief gratings}, JOSA A 14, 2758 (1997).
  \item E.~Popov, M.~Nevière, \enquote{Maxwell equations in Fourier space: fast-converging formulation for diffraction by arbitrary shaped, periodic, anisotropic media}, JOSA A 18, 2886 (2001).
  \item K.~Watanabe, R.~Petit, M.~Nevière, \enquote{Differential theory of gratings made of anisotropic materials}, JOSA A 19, 325 (2002).
  \item T.~Schuster, J.~Ruoff, N.~Kerwien, S.~Rafler, W.~Osten, \enquote{Normal vector method for convergence improvement using the RCWA for crossed gratings}, JOSA A 24, 2880 (2007).
  \item P.~Götz, T.~Schuster, K.~Frenner, S.~Rafler, W.~Osten, \enquote{Normal vector method for the RCWA with automated vector field generation}, Opt.\ Express 16, 17295 (2008).
  \item E.~Hairer, G.~Wanner, \emph{Solving Ordinary Differential Equations II}, Springer (Butcher-Tableaus Gauß/Radau).
  \item NVIDIA cuSOLVER Documentation, \verb|cusolverDnXgeev| (seit CUDA 12.6 Update 2).
\end{enumerate}
(Literaturangaben 7, 8 vor Verwendung bibliographisch prüfen.)
